\section{Related Work}
\label{sec:related}

The central thesis of this paper is that the analogy framework of \textit{computer architecture} can be used to envision the full layered design of future model-native computing systems.
This thesis did not emerge in a vacuum---by 2025, a substantial body of work had independently touched upon key insights such as ``LLM systems need OS-style abstractions,'' ``agents require memory management,'' and ``tool invocation demands standardized protocols.''
This section surveys the most relevant literature along functional themes, progressively revealing the contribution and limitation of each line of work, and ultimately positioning the unique contribution of this paper.

\subsection{LLM-as-OS / Agent Kernel}
\label{sec:related-llmos}

The earliest systematic articulation of the ``LLM as operating system'' analogy is due to Ge et al.\ (2023), who proposed ``LLM as OS, Agents as Apps''~\cite{ge2023llmos}.
Their central insight can be summarized by a system mapping diagram: the LLM corresponds to the OS kernel (responsible for core inference and scheduling), the \textit{context window} corresponds to main memory, external storage corresponds to the file system, hardware tools correspond to peripherals, software tools correspond to programming libraries, and user prompts correspond to user commands.
In other words, they reconceived the LLM not as a ``chatbot'' but as a ``kernel process''---a metaphor later popularized by Karpathy on social media~\cite{karpathy2023llmos}.
Within the analogy framework of this paper, the work of Ge et al.\ directly corresponds to the embryonic notion of the Layer~2 Agent OS.

Mei et al.'s AIOS~\cite{mei2024aios} advanced this line of thinking from an analogy to a runnable system implementation.
The core design of AIOS decouples responsibilities that were previously entangled with application logic in the agent runtime and places them into a dedicated ``kernel,'' including:
(1)~a scheduler managing concurrent execution of multiple agents;
(2)~context management distributing the limited context window across agents;
(3)~memory management handling reads and writes of short-term and long-term memory;
(4)~storage management governing files and persistent data; and
(5)~access control governing agent permissions over tools and data.
Experimental results show that this decoupled design yields up to $2.1\times$ speedup.
In our analogy framework, AIOS realizes an early Agent OS prototype whose kernel responsibilities---process management, memory management, and file system modules---map directly onto those of a traditional operating system.

AgentOS~\cite{agentos2025} went further by defining the LLM as a \textit{reasoning kernel} constrained by ``structured operating system logic''---that is, the LLM is not a free-form inference engine but a controlled reasoning core that must operate within the structural framework of OS logic.
This perspective emphasizes that the LLM's reasoning capabilities must be orchestrated and constrained by structured system logic, rather than allowing the model to generate without restraint.

By 2025--2026, this line of research advanced from academic prototypes toward real desktop and mobile scenarios.
Microsoft's UFO$^2$~\cite{ufo2_2025} organizes a HostAgent (for understanding user intent), an AppAgent (for operating specific applications), a GUI--API hybrid action layer, and virtual-desktop parallel operations into a ``Desktop AgentOS.''
Its innovation lies in simultaneously supporting both GUI screenshot comprehension and native API calls as interaction modalities.
Aura~\cite{aura2025}, targeting mobile scenarios, adopts a hub-and-spoke topology: a privileged System Agent parses user intent, while multiple sandboxed App Agents each handle a specific domain.
Its design philosophy is ``security first''---inter-agent communication must be mediated through the System Agent, analogous to how inter-process communication in a traditional OS must pass through the kernel.

\textbf{Relationship to this paper.}
The works above each define a role for the LLM within the system---OS kernel, reasoning kernel, or Agent Kernel---but they have not converged on a unified layered model.
The ICAM model proposed in Section~\ref{sec:icam} consolidates these disparate ``kernel'' concepts into a six-layer architecture and explicitly positions the LLM as Layer~1 (the probabilistic compute core) while the Agent OS occupies Layer~2 as an independent abstraction---the two should not be conflated.

\subsection{Memory Hierarchy and Virtual Context Management}
\label{sec:related-memory}

Large language models have a finite \textit{context window}---for example, GPT-4 Turbo supports 128K tokens, Claude~3 supports 200K tokens, and Gemini~1.5~Pro reaches up to 1M tokens~\cite{gemini2024}.
Yet regardless of how large the model's window grows, the usable capacity remains bounded, which is structurally identical to the ``finite physical memory'' problem in classical computing.
A series of works has therefore borrowed \textbf{memory management} techniques from operating systems to expand the effective context available to agents.

MemGPT~\cite{memgpt2023} (later evolved into the Letta platform) is the foundational work in this direction.
Rather than simply attaching a retrieval-augmented generation (RAG)~\cite{lewis2020rag} pipeline to look up external knowledge, MemGPT draws directly on the \textbf{virtual memory management} mechanism of a traditional OS:
it treats the limited context window as ``main memory,'' external databases as ``disk,'' and employs an OS-like memory manager that automatically pages information fragments in and out between the context window and external storage, giving the agent the illusion of unbounded context.
MemGPT's control flow also explicitly employs \textit{interrupts}---when information must be retrieved from external storage, an interrupt transfers control to the memory manager.
This design maps directly onto the ``virtual memory'' concept in our analogy.

MemOS~\cite{memos2025} elevates the problem to a more systematic plane.
It unifies three granularities of memory into manageable system resources:
(1)~\textit{plaintext memory}---external knowledge stored as text;
(2)~\textit{activation-based memory}---internal hidden-layer representations of the model (analogous to CPU registers); and
(3)~\textit{parameter-level memory}---knowledge encoded in model parameters (analogous to firmware/ROM).
MemOS uses the ``MemCube'' as an encapsulation unit, bundling content with provenance and versioning metadata, and attempts to build a migration bridge between retrieval and parameter learning.
This work corresponds to the ``multi-level storage hierarchy from KV cache to parameter storage'' in our analogy.

MemoryOS~\cite{memoryos2025} applies OS memory management principles directly to the agent memory system.
HiAgent~\cite{hiagent2024} (published at ACL~2025) proposes ``hierarchical working memory management''---using subgoals as memory chunking units to organize long-horizon task memory into a hierarchy, analogous to multi-level page tables in an operating system.
A-MEM~\cite{amem2025} introduces the concept of \textit{agentic memory}: memory is not a passively stored record but a dynamic resource that agents can autonomously organize, recombine, and update.
LongMemEval~\cite{longmemeval2024} further demonstrates that long-term memory is not simply ``storing more chat logs,'' but a complex systems engineering challenge involving information extraction, temporal reasoning, knowledge updating, and the ability to refuse to answer.

\textbf{Relationship to this paper.}
Research on memory subsystems has become highly systematic, yet it is typically discussed in isolation from the agent runtime and I/O layers, lacking a full-stack perspective.
More importantly, existing work has not explicitly identified the \textit{KV cache} (the key-value cache maintained during model inference) as the counterpart of ``hardware cache'' in the analogy---whereas this paper argues that the KV cache is precisely the key hardware abstraction bridging the probabilistic compute core and the semantic memory hierarchy.
We elaborate on this correspondence in Section~\ref{sec:kv}.

\subsection{Agent Frameworks and Runtimes}
\label{sec:related-agent-framework}

If the preceding two themes define the ``kernel'' and ``memory,'' respectively, this section focuses on how programs execute on top of that kernel---that is, the execution frameworks and runtime environments for agents.

\textbf{Foundational capabilities: interleaving reasoning and action.}
ReAct~\cite{yao2023react} is the foundational work on agent execution paradigms.
Prior to ReAct, reasoning (e.g., Chain-of-Thought) and acting (e.g., invoking tools) were separate, sequential steps for LLMs.
ReAct demonstrated that \textbf{interleaving} reasoning and action within a single loop---``think one step, act one step, observe the result, think again''---significantly improves task completion quality.
This is analogous to the fetch-decode-execute cycle in traditional computing, where the CPU executes an instruction, accesses memory, and checks a condition in tight alternation.

Toolformer~\cite{schick2023toolformer} addressed another foundational question: enabling the model to learn \textbf{when} to call APIs on its own.
By automatically inserting API call examples into training data, Toolformer empowers the model to autonomously decide during generation whether it needs to invoke a particular tool---analogous to the \textit{system call} mechanism in an operating system, where a user program initiates a request to the kernel.

\textbf{The rise of coding agents.}
The most striking agent advances in 2024--2026 have concentrated in software engineering.
OpenAI's Codex~CLI~\cite{openai_codex_overview_2026}, released in April~2025, is a terminal-based agent programming assistant.
It can read code repositories, edit files, execute commands, and operate safely within a sandboxed environment~\cite{openai_codex_sandbox_2026}.
In early 2026, Codex~CLI introduced a subagent feature~\cite{openai_codex_subagents_2026}---the main agent can spawn specialized subagents that execute subtasks in parallel, each with its own independent context window and sandbox.
This directly mirrors the \texttt{fork()} model in traditional operating systems, where a parent process creates child processes.
Codex also supports custom instructions via AGENTS.md files~\cite{openai_codex_agents_md_2026} and reusable capability modules through Skills~\cite{openai_codex_skills_2026}.

Anthropic's Claude Code~\cite{anthropic_claudecode_overview_2026} is another representative agent runtime.
Like Codex, it is a terminal-native coding agent capable of editing files, executing shell commands, and connecting to external services.
An independent research study systematically analyzed Claude Code's system architecture~\cite{dive_claudecode_2026} and found that only approximately 1.6\% of its codebase constitutes AI decision logic, while the remaining 98.4\% is \textit{operational harness}---including safety controls, tool orchestration, and context management.
Claude Code likewise supports subagents~\cite{anthropic_claudecode_subagents_2026} and Skills-based extensions~\cite{anthropic_claudecode_skills_2026}, and provides project-level memory through CLAUDE.md files~\cite{anthropic_claudecode_memory_2026}.

Notably, the temporal span of agent execution is expanding from minutes to hours or even days.
Rakuten engineers tested Claude Code on the vLLM project (12.5~million lines of code) to perform activation vector extraction; the agent ran autonomously for 7~hours and achieved 99.9\% numerical accuracy~\cite{anthropic2026agenticreport}.
This trend means that agent runtimes must handle state persistence, checkpoint recovery, and cross-step resource management far beyond what current interactive sessions require.

OpenHands~\cite{openhands_paper_2025} (formerly OpenDevin) is an open-source general-purpose agent platform published at ICLR~2025.
It achieves approximately 70--74\% issue-resolution rates on the SWE-bench Verified benchmark~\cite{swebench2023} and supports multiple LLM backends.
SWE-agent~\cite{sweagent2024} focuses specifically on software engineering tasks, achieving comparable performance with a more streamlined Agent-Computer Interface (ACI) design.

\textbf{Relationship to this paper.}
These agent frameworks already exhibit operating-system-like characteristics in their architecture: Codex's subagents correspond to process management, Claude Code's CLAUDE.md corresponds to configuration and environment management, and OpenHands' ACI corresponds to the system call interface.
However, they remain independently designed ``application-layer'' tools lacking a unified layered abstraction.
The ICAM model proposed in this paper provides a unifying theoretical framework that maps their functionality onto well-defined architectural layers.

\subsection{Tool Calling and I/O Protocols (MCP, A2A)}
\label{sec:related-io}

Agents must interact with the external world---querying databases, calling APIs, manipulating file systems.
In classical computing, this requirement corresponds to the I/O subsystem (input/output devices and buses).
The key development over the past two years is that I/O abstractions are moving from ad hoc tool-calling conventions toward \textbf{standardized protocols}.

HuggingGPT~\cite{hugginggpt2023} is an early representative work: it places the LLM in a ``controller'' position, tasking it with parsing user requests, selecting appropriate expert models, coordinating execution, and aggregating results.
This is analogous to the device driver management layer in an operating system, where the OS uniformly manages calls to different devices (models).
Gorilla~\cite{patil2024gorilla} focuses on a more specific problem: when API documentation is frequently updated, how does the model maintain call accuracy? This corresponds to the version-compatibility problem for device drivers.

\textbf{MCP: the model-to-system I/O bus.}
In November~2024, Anthropic introduced the Model Context Protocol (MCP)~\cite{mcp_intro_2026}, defined as an open standard for connecting AI applications with external data sources, tools, and workflows.
MCP adopts a client--server architecture: an MCP Client runs on the AI application side, while an MCP Server encapsulates a specific tool or data source.
Through a unified JSON-RPC protocol, any MCP-compatible AI application can connect to any tool that provides an MCP Server---Anthropic has likened MCP to ``USB-C for AI.''
By 2026, the MCP ecosystem has expanded rapidly, with Claude Code~\cite{anthropic_claudecode_mcp_2026} and numerous third-party tools all supporting MCP connections.

\textbf{A2A: the agent-to-agent interconnect protocol.}
In April~2025, Google announced the Agent-to-Agent Protocol (A2A)~\cite{google_a2a_2025} at Cloud Next~'25, explicitly positioned as a complementary open protocol to MCP.
A2A addresses not the ``model-to-tool'' connection problem but the ``agent-to-agent'' collaboration problem---supporting capability discovery, task lifecycle management, long-running tasks, and multimodal coordination.
A2A was donated to the Linux Foundation in 2025~\cite{a2alinux2025}, becoming an industry standard candidate.
A2A employs a peer-to-peer client--server model in which every agent can act as both a requesting client and a responding server.

\textbf{Relationship to this paper.}
Taken together, MCP plays the role of a model-to-system I/O bus (analogous to USB/PCIe), while A2A serves as an agent-to-agent interconnect protocol (analogous to the TCP/IP networking stack).
In the ICAM model proposed in this paper, these correspond to Layer~4 (I/O and tool layer) and Layer~5 (interconnect and protocol layer), respectively.
No existing work has yet integrated both protocols into a unified layered abstraction or examined their interactions with context management, kernel scheduling, and other layers at the architectural level.

\subsection{Multi-Agent Collaboration as Distributed Systems}
\label{sec:related-multiagent}

When multiple agents collaborate, the system exhibits characteristics of distributed computing---a development that parallels the historical evolution from standalone machines to distributed systems.

AutoGen~\cite{wu2023autogen} (Microsoft Research, 2023) is the foundational framework for multi-agent systems.
Its core design treats \textit{multi-agent conversation} as universal infrastructure: agents in different roles coordinate tasks through natural language dialogue.
AutoGen has evolved to v0.4 and has been integrated into the Microsoft Agent Framework, supporting scalable, production-grade multi-agent systems.

MetaGPT~\cite{metagpt2023} (ICLR~2024 Oral) introduced a critical design idea: explicitly encoding the \textbf{Standard Operating Procedure} (SOP) from human organizations into multi-agent collaboration.
Specifically, MetaGPT simulates a virtual software company: a product manager writes requirements, an architect designs the system, engineers write code, and a QA engineer tests---each role is an agent, and the SOP defines the workflow and information flow among them.
This is analogous to \textit{workflow orchestration} in distributed systems, constraining unordered agent interactions into an ordered assembly line.

This pipeline pattern has been validated at industrial scale.
Fountain's Copilot system employs a central orchestrator agent that coordinates three specialized sub-agents for candidate screening, document generation, and sentiment analysis.
This hierarchical multi-agent deployment achieved 50\% faster candidate screening, 40\% faster onboarding, and $2\times$ candidate conversion rates, compressing the full staffing center configuration cycle from over one week to under 72~hours~\cite{anthropic2026agenticreport}.
This hub-and-spoke topology contrasts sharply with MetaGPT's pipeline topology and AutoGen's peer-to-peer topology, suggesting that coordination topology itself is a key design variable influencing multi-agent system performance~\cite{wu2023autogen,metagpt2023}.

The Internet of Agents (IoA)~\cite{ioa2024} (published at ICLR~2025) envisions the agent ecosystem from an ``internet'' rather than single-system perspective.
Its core contribution is an agent integration protocol supporting dynamic team formation, collaboration, and distributed execution among heterogeneous agents.
IoA draws design inspiration directly from the Internet: just as the Internet connects computers from different manufacturers running different operating systems, IoA aims to connect agents built with different frameworks by different organizations.
This has also spurred IETF-level standardization efforts (IoA Protocol Internet Draft).

\textbf{Relationship to this paper.}
Multi-agent systems already exhibit ``distributed system'' morphology at the object-of-study level: AutoGen corresponds to message-passing mechanisms, MetaGPT to workflow orchestration, and IoA to network interconnect protocols.
However, most work remains at the ``framework and protocol'' layer without integrating with individual agent memory management, kernel scheduling, I/O, and access control into a unified architecture.
The ICAM model maps multi-agent collaboration onto Layer~5 (interconnect and protocol layer) and Layer~6 (application layer), and discusses cross-layer interaction design principles.

\subsection{Cognitive Architectures and Security Governance}
\label{sec:related-cognition-security}

The preceding subsections have surveyed related work from a functional systems perspective (kernel, memory, runtime, I/O, multi-agent).
This section turns to two higher-level themes: what cognitive architecture should agent systems follow, and how should their security be governed?

\textbf{Cognitive architectures.}
CoALA (Cognitive Architectures for Language Agents)~\cite{sumers2024coala} provides a theoretical framework grounded in cognitive science.
CoALA organizes language agents into three core components:
(1)~\textit{modular memory}---divided into short-term working memory and long-term memory (semantic, episodic, and procedural);
(2)~\textit{structured action space}---internal actions (reasoning, memory retrieval) and external actions (tool invocation, environment interaction); and
(3)~a \textit{generalized decision-making process}---a plan--execute--observe loop.
The value of CoALA lies in providing the theoretical anchor for ``why this system is not merely a software stack assembly, but a genuine cognitive architecture.''
L2MAC~\cite{holt2024l2mac} (ICLR~2024) explicitly frames its architecture as the ``first practical LLM-based stored-program automatic computer (von Neumann architecture),'' comprising an instruction register, file storage, control unit, and independent LLM agent modules.
Mi et al.~\cite{mi2025llmagentscomputersystems} organize their LLM agent survey from a ``computer systems perspective,'' explicitly inspired by the von Neumann architecture: they decompose agents into perception, cognition, memory, tool, and action modules, observe that the context window is analogous to main memory and databases to disk, and point out that agents currently lack a ``cache module.''

\textbf{Security governance.}
ArbiterOS~\cite{arbiteros2025} identifies a fundamental problem: we are using the mental model of \textbf{deterministic software engineering} to direct what are \textbf{inherently probabilistic processors}.
This contradiction demands the introduction of a governance layer at the system architecture level.
ArbiterOS accordingly proposes the ``Agentic Computer'' mental model---redefining the LLM as a \textit{Probabilistic CPU} and designing governance-layer concepts including a governor, a formal instruction set, and a hardware abstraction layer (HAL).
Within our analogy framework, ArbiterOS's core insight---that models are probabilistic and require a deterministic governance layer---is one of the direct inspirations for the \textit{dual-plane architecture} (probabilistic execution plane + deterministic control plane) proposed in this paper.

CaMeL~\cite{debenedetti2025camel} (Google DeepMind, 2025) applies operating-system capability-security principles to LLM agents from a safety perspective.
Its core design philosophy is ``Don't execute data''---strictly separating trusted instructions from untrusted data, extracting control flow and data flow and enforcing separate security policies on each to defend against prompt injection attacks.
IronClaw~\cite{ironclaw2025} (2026) explicitly advocates protecting AI agents in the same way an operating system protects its processes---migrating OS security mechanisms such as privilege separation, sandbox isolation, and audit logging into the agent system.

\textbf{Relationship to this paper.}
The cognitive architecture line provides theoretical foundations, while the security governance line supplies engineering constraints.
Yet these two threads remain unintegrated---cognitive architectures rarely address security, and security governance rarely considers cognitive structure.
The dual-plane architecture proposed in this paper attempts to unify both: the probabilistic execution plane carries cognitive capabilities, while the deterministic control plane enforces security governance, and the two interact through well-defined interface contracts.

\subsection{Common Boundaries and Unsolved Problems}
\label{sec:related-gap}

The six themes surveyed above collectively approach the vision of a ``model-native computing architecture,'' yet four critical problems remain unsolved in a unified manner.

\textbf{Problem 1: What role does the LLM actually play?}
Existing work exhibits a fundamental divergence in positioning the LLM within the system.
Ge et al.\ and AIOS place the LLM at the core of the OS~\cite{ge2023llmos,mei2024aios}---the LLM \textit{is} the kernel;
ArbiterOS demotes it to a ``Probabilistic CPU'' governed by the kernel~\cite{arbiteros2025}---the LLM is a managed processor;
AgentOS defines it as a reasoning kernel constrained by structured OS logic~\cite{agentos2025}---the LLM is a controlled reasoning core.
These three positions lead to fundamentally different system designs.
The ICAM model proposed in this paper offers a unified answer: the LLM is Layer~1 (probabilistic compute core / CPU), and the Agent OS is Layer~2 (operating system kernel); they are distinct abstraction layers and should not be conflated.
The LLM provides compute capability, while the Agent OS provides management and scheduling---just as the CPU and the OS kernel are two independent yet tightly cooperating layers in classical computer architecture.

\textbf{Problem 2: No full-stack layered model exists.}
No mainstream work simultaneously places the probabilistic compute core, semantic memory hierarchy, tool and peripheral bus, Agent OS/kernel, interconnect protocols, governance and security layer, programming model and instruction semantics, and evaluation metrics within a single architectural framework.
L2MAC~\cite{holt2024l2mac} focuses on the stored-program computer model; Mi et al.~\cite{mi2025llmagentscomputersystems} provide a von Neumann-inspired survey framework, but each covers only one facet of the architecture.
Memory management work does not discuss I/O protocols; agent frameworks do not discuss KV cache; security governance does not discuss multi-agent collaboration---each line of work makes genuine contributions, but their boundaries are clearly delineated.

\textbf{Problem 3: The probabilistic--deterministic tension has not been codified architecturally.}
ArbiterOS~\cite{arbiteros2025} correctly identifies the tension between probabilistic processors and deterministic governance, but it proposes only a conceptual ``governance-first'' paradigm.
CaMeL~\cite{debenedetti2025camel} and IronClaw~\cite{ironclaw2025} provide concrete security mechanisms, but they are not embedded within a complete layered architecture.
This paper \textit{codifies this tension architecturally} through the dual-plane architecture: the probabilistic execution plane carries the LLM's inference capabilities, while the deterministic control plane enforces OS-level management and security responsibilities, with the two planes cooperating through explicit interface protocols.

These first three unsolved problems define the contribution space of this paper: providing a unified, layered, dual-plane architectural model that integrates all six threads within a single master diagram.

\textbf{Problem 4: Structural constraints of human--agent collaboration have not been modeled.}
Existing work treats the human as a boundary condition of the system (Layer~6 user), but the latest industrial evidence reveals a ``collaboration paradox'': engineers use AI for approximately 60\% of their work, yet can fully delegate only 0--20\% of their tasks~\cite{anthropic2026agenticreport}.
Engineers tend to delegate tasks that are ``easy to verify or low-risk'' while retaining those that are ``conceptually difficult or depend on design decisions.''
This structural constraint suggests that the human--agent interaction loop should become an explicit component of the architecture, rather than an external assumption.
